\clearpage
\hypertarget{abhaken}{}\einheitenkopf{Das kann ich}
\zweigzeile{Hake ab, was du kannst.}
\par\medskip\noindent\begin{minipage}{\linewidth}
\textbf{Einheit 1 · Proportionale Funktion}\par\smallskip
\noindent\begin{tabular}{@{}p{7mm} p{14cm} c@{}}
\hfill 1. & Ich kann die Gerade zu $y = m \cdot x$ zeichnen. & $\square$ \\
\hfill 2. & Ich kann die Gerade zu $y = m \cdot x$ zeichnen – weiter. & $\square$ \\
\hfill 3. & Ich kann Werte an einer Ursprungsgeraden ablesen und berechnen. & $\square$ \\
\hfill 4. & Ich kann erkennen, ob eine proportionale Funktion vorliegt, und ihre Gleichung angeben. & $\square$ \\
\hfill 5. & Ich finde den Fehler beim Bestimmen des Faktors. & $\square$ \\
\hfill 6. & Ich kann Eigenschaften von Ursprungsgeraden begründen. & $\square$ \\
\hfill 7. & Ich kann eine Sachsituation mit $y = m \cdot x$ beschreiben. & $\square$ \\
\end{tabular}
\end{minipage}
\par\medskip\noindent\begin{minipage}{\linewidth}
\textbf{Einheit 2 · Lineare Funktion $f(x) = m \cdot x + n$}\par\smallskip
\noindent\begin{tabular}{@{}p{7mm} p{14cm} c@{}}
\hfill 8. & Ich erkenne $m$ und $n$ in der Gleichung. & $\square$ \\
\hfill 9. & Ich erkenne, ob eine Gerade steigt oder fällt. & $\square$ \\
\hfill 10. & Ich kann ein Steigungsdreieck lesen. & $\square$ \\
\hfill 11. & Ich kann eine Gerade mit einer Wertetabelle zeichnen. & $\square$ \\
\hfill 12. & Ich kann eine Gerade mit $n$ und Steigungsdreieck zeichnen. & $\square$ \\
\hfill 13. & Ich kann eine Gerade mit $n$ und Steigungsdreieck zeichnen – weiter. & $\square$ \\
\hfill 14. & Ich kann $n$ und $m$ an einer Geraden ablesen. & $\square$ \\
\hfill 15. & Ich kann einer Gleichung die passende Gerade zuordnen. & $\square$ \\
\hfill 16. & Ich kann an $m$ und $n$ ablesen, wie Geraden zueinander liegen. & $\square$ \\
\hfill 17. & Ich finde den Fehler beim Ablesen der Steigung. & $\square$ \\
\hfill 18. & Ich kann begründen, welche Gerade steiler ist. & $\square$ \\
\hfill 19. & Ich kann $m$ und $n$ in einer Sachsituation deuten. & $\square$ \\
\end{tabular}
\end{minipage}
\par\medskip\noindent\begin{minipage}{\linewidth}
\textbf{Einheit 3 · Punkte und Werte}\par\smallskip
\noindent\begin{tabular}{@{}p{7mm} p{14cm} c@{}}
\hfill 20. & Ich kann eine Zahl für $x$ einsetzen. & $\square$ \\
\hfill 21. & Ich kann Funktionswerte berechnen. & $\square$ \\
\hfill 22. & Ich kann zu einem Funktionswert das passende $x$ berechnen. & $\square$ \\
\hfill 23. & Ich kann die Nullstelle einer linearen Funktion berechnen. & $\square$ \\
\hfill 24. & Ich kann die Nullstelle einer linearen Funktion berechnen – weiter. & $\square$ \\
\hfill 25. & Ich kann prüfen, ob ein Punkt auf einer Geraden liegt. & $\square$ \\
\hfill 26. & Ich kann eine Wertetabelle ausfüllen und einer Funktion zuordnen. & $\square$ \\
\hfill 27. & Ich kann Nullstellen und Schnittpunkte am Graphen ablesen. & $\square$ \\
\hfill 28. & Ich finde den Fehler bei der Nullstelle. & $\square$ \\
\hfill 29. & Ich kann Aussagen über eine Gerade beurteilen. & $\square$ \\
\hfill 30. & Ich kann begründen, wie viele Nullstellen eine Gerade hat. & $\square$ \\
\hfill 31. & Ich kann Funktionswert und Nullstelle im Sachzusammenhang deuten. & $\square$ \\
\end{tabular}
\end{minipage}
\par\medskip\noindent\begin{minipage}{\linewidth}
\textbf{Einheit 4 · Gleichung bestimmen}\par\smallskip
\noindent\begin{tabular}{@{}p{7mm} p{14cm} c@{}}
\hfill 32. & Ich kann die Gleichung einer Geraden am Graphen ablesen. & $\square$ \\
\hfill 33. & Ich kann die Gleichung aus der Steigung und einem Punkt bestimmen. & $\square$ \\
\hfill 34. & Ich kann die Gleichung einer Geraden durch zwei Punkte bestimmen. & $\square$ \\
\hfill 35. & Ich kann eine Gerade durch zwei Punkte zeichnen. & $\square$ \\
\hfill 36. & Ich kann den Schnittpunkt zweier Geraden berechnen. & $\square$ \\
\hfill 37. & Ich kann den Schnittpunkt zweier Geraden berechnen – weiter. & $\square$ \\
\hfill 38. & Ich finde den Fehler bei der Steigung aus zwei Punkten. & $\square$ \\
\hfill 39. & Ich kann begründen, wie zwei Geraden zueinander liegen. & $\square$ \\
\hfill 40. & Ich kann eine Geradengleichung im Sachzusammenhang bestimmen. & $\square$ \\
\end{tabular}
\end{minipage}
\par\medskip\noindent\begin{minipage}{\linewidth}
\textbf{Einheit 5 · Anwendungen}\par\smallskip
\noindent\begin{tabular}{@{}p{7mm} p{14cm} c@{}}
\hfill 41. & Ich finde Anfangswert und Änderung im Text. & $\square$ \\
\hfill 42. & Ich kann die passende Gleichung zu einem Tarif erkennen. & $\square$ \\
\hfill 43. & Ich kann eine Gleichung zu einer Sachsituation aufstellen. & $\square$ \\
\hfill 44. & Ich kann berechnen, welcher Wert nach einer bestimmten Zeit erreicht ist. & $\square$ \\
\hfill 45. & Ich kann zwei Tarife für eine bestimmte Nutzung vergleichen. & $\square$ \\
\hfill 46. & Ich kann berechnen, ab wann ein Tarif günstiger ist. & $\square$ \\
\hfill 47. & Ich finde den Fehler beim Aufstellen einer Tarifgleichung. & $\square$ \\
\hfill 48. & Ich kann einem Tarif seinen Graphen zuordnen. & $\square$ \\
\hfill 49. & Ich kann zu einer Gleichung eine passende Situation beschreiben. & $\square$ \\
\end{tabular}
\end{minipage}
